\section*{Data Management Plan}
\addcontentsline{toc}{section}{\protect\numberline{}Data Management Plan}

Due to the nature of the projects, the data used in the projects will mainly be synthetic data generated from simulations of math and statistics models.
PI has had extensive experience gathering, archiving with large quantities of data which contain sensitive information. 

% PI and researchers to be supported by the grant have had human subjects research through the Collaborative Institutional Training Initiative (CITI) program. They have also completed Responsible Conduct of Research Courses which details data management requirements of human subjects data. PI has also had extensive publication and project experience in the areas of human-computer interaction, machine learning, social computing and other areas of conducting user studies. Researchers and students supported by the project are also expected to complete all HSR training through CITI which detail how to properly manage, sharing, and analyze data containing private data. Additionally, PI has had extensive engineering experience in industry developing and maintaining enterprise-level data management systems. This expertise will help to ensure the secure access and backup of data. 

\subsection*{Data Sharing}
% NA
% Georgia State University provides faculty with 100GB of cloud storage as well as 1TB for students through OneDrive. This will be used to share files between researchers. T
The PI will administer the access of this data. 
% Additionally, the PI's lab hosts a Network Attached Storage (NAS) system by Synology which can be configured to share files over the internet with up to 25Gbps. The NAS device currently has over 30TB of total storage and can be expanded to 140TB of data if necessary. This will host any large files that need to be stored and shared between collaborators at each institution. The NAS is configured to have a two-drive failure tolerance which will ensure that no data is lost. Access to all cloud and NAS storage devices are password protected with strong password protection and will be managed by the PI. 
% 
Results and data from the project will be submitted to reputable journals and ML conferences such as SIAM uncertainty quantification, Annals of Statistics, and Neurips. We expect our work to be presented at a variety of workshops, conferences, and journals due to its high impact. Properly anonymized data will be shared through data sharing platforms such as the OpenScience Framework. We will also utilize our NAS or Dropbox link to share data which may be too large. GitHub will also be used to share software and data for public use. 
\subsection*{Handling Sensitive Data}
% NA
% Data gathered from human subjects will not contain personally identifiable information such as their names, government issued identifiers, or others. A unique randomly generated ID will be assigned to each subject and does not map to any personal information. Additionally, this data gathered will have the permission of any participants for metric/questionnaires. Data will be stored on servers hosted at our respective institutions which are typically only accessible through when connected to the institutional network using a Virtual Private Network (VPN). Additionally, servers will be password protected so that login via secure shell (SSH) or other mechanism is required for accessing any stored data. Servers also have strong password requirements to protect against insecure access. Network access is also over an encrypted and secure network to protect against any attacks. 
Data will be stored on researcher’s password-protected laptops.
% Any data which is shared will adhere to the strict processes laid out in Human Subjects Research requirements to prevent leaking of sensitive information. 
The PI on the project will ensure that data is handled and disseminated in a privacy-protected, safe, and responsible manner. 
\subsection*{Authorization for data access and protection of data}
% NA
Authorization will be monitored and supported by the PI. This can be configured to provide password protected access to specific files. We will primarily use this system as a storage device for data during research stages. Once the data is to be shared, we will utilize an approval system to share if necessitated by the Internal Review Board regarding the sharing of collected data. If additional monitoring of shared data is necessary, a web form will be set up with terms of use to ensure that other users of the data properly protect the identity of individuals.
